\section{Equivalencia matemática entre modelos de difusión y modelos de flujo}

Esta sección es la parte teórica central de esta clase, donde el ponente deriva detalladamente la relación matemática exacta entre los modelos de difusión y los modelos de flujo.

\subsection{Partiendo del modelo de difusión variacional}

Recuerde que en el Modelo de Difusión Variacional (Variational Diffusion Model, VDM), el proceso hacia adelante se define como:
$$q(\mathbf{x}_t \mid \mathbf{x}_0) = \mathcal{N}(\mathbf{x}_t; \alpha_t \mathbf{x}_0, \sigma_t^2 \mathbf{I})$$

En el contexto de Flow Matching, interprete este proceso hacia adelante nuevamente como un mapeo de flujo condicional:
$$\psi_t(\mathbf{x}_0 \mid \mathbf{x}_1) = \sigma_t \cdot \mathbf{x}_0 + \alpha_t \cdot \mathbf{x}_1$$

\begin{importantbox}{Insight central: los modelos de difusión son un caso especial de modelos de flujo}
Cuando el mapeo de flujo condicional se define en la forma gaussiana lineal anterior y las funciones $\alpha_t, \sigma_t$ coinciden exactamente con las funciones de programación del modelo de difusión variacional, el modelo de difusión puede considerarse un\textbf{caso especial} de los modelos de flujo. En otras palabras, los modelos de flujo son una\textbf{generalización} de los modelos de difusión. La única diferencia entre ambos radica en:
\begin{itemize}
    \item \textbf{Modelos de difusión}: primero se define la trayectoria de probabilidad (proceso de adición de ruido hacia adelante), luego se derivan el campo vectorial y las trayectorias
    \item \textbf{Modelos de flujo}: primero se define el mapeo de flujo condicional (trayectoria), luego se derivan el campo vectorial y la trayectoria de probabilidad
\end{itemize}
\end{importantbox}

\subsection{Reaparecimiento de la fórmula de Tweedie}

Aquí, el ponente introduce una herramienta de cálculo importante: la fórmula de Tweedie. Dado un punto de datos intermedio $\mathbf{x}_t = \sigma_t \mathbf{x}_0 + \alpha_t \mathbf{x}_1$, donde $\mathbf{x}_0 \sim \mathcal{N}(\mathbf{0}, \mathbf{I})$ y $\mathbf{x}_1 \sim p_{\text{data}}$, podemos calcular la esperanza a posteriori de $\mathbf{x}_1$:

$$\mathbb{E}[\mathbf{x}_1 \mid \mathbf{x}_t] = \frac{\mathbf{x}_t - \sigma_t \cdot \mathbb{E}[\mathbf{x}_0 \mid \mathbf{x}_t]}{\alpha_t} = \frac{\mathbf{x}_t + \sigma_t^2 \nabla_{\mathbf{x}_t} \log p_t(\mathbf{x}_t)}{\alpha_t}$$

Donde:
\begin{itemize}
    \item $\nabla_{\mathbf{x}_t} \log p_t(\mathbf{x}_t)$ es la función de puntaje (Score Function)
    \item Esto coincide exactamente con la fórmula de Tweedie presentada en la Clase 5, solo que se utilizan las convenciones de notación de Flow Matching
\end{itemize}

\begin{knowledgebox}{Intuición detrás de la fórmula de Tweedie}
La idea central de la fórmula de Tweedie es: dado una observación $\mathbf{x}_t$ contaminada con ruido, la media a posteriori de la señal original $\mathbf{x}_1$ puede calcularse mediante la función de puntaje. La función de puntaje apunta en la dirección donde crece más rápido la densidad de probabilidad del logaritmo, es decir, la dirección de ``mayor concentración de datos''. En un modelo de verosimilitud gaussiana, esto equivale a moverse desde la posición actual hacia la dirección de la media a posteriori.
\end{knowledgebox}

\subsection{Relación entre el campo de velocidad y la función de puntaje}

Esta es una de las derivaciones más importantes de esta clase. Partiendo del campo vectorial condicional:

$$u_t(\mathbf{x} \mid \mathbf{x}_1) = \dot{\sigma}_t \cdot \mathbf{x}_0 + \dot{\alpha}_t \cdot \mathbf{x}_1$$

Expresando $\mathbf{x}_0$ en términos de $\mathbf{x}_t$ y $\mathbf{x}_1$: $\mathbf{x}_0 = \frac{\mathbf{x}_t - \alpha_t \mathbf{x}_1}{\sigma_t}$, y sustituyendo obtenemos:

$$u_t(\mathbf{x}_t \mid \mathbf{x}_1) = \frac{\dot{\sigma}_t}{\sigma_t}(\mathbf{x}_t - \alpha_t \mathbf{x}_1) + \dot{\alpha}_t \mathbf{x}_1$$

$$= \frac{\dot{\sigma}_t}{\sigma_t} \mathbf{x}_t + \left(\dot{\alpha}_t - \frac{\dot{\sigma}_t \alpha_t}{\sigma_t}\right) \mathbf{x}_1$$

Tomando la esperanza sobre $\mathbf{x}_1$ (marginalizando) y utilizando la fórmula de Tweedie, obtenemos el campo vectorial marginal:

\begin{importantbox}{Fórmula equivalente entre el campo de velocidad y la función de puntaje}
$$u_t(\mathbf{x}_t) = \frac{\dot{\sigma}_t}{\sigma_t} \mathbf{x}_t + \left(\dot{\alpha}_t - \frac{\dot{\sigma}_t \alpha_t}{\sigma_t}\right) \mathbb{E}[\mathbf{x}_1 \mid \mathbf{x}_t]$$

Sustituyendo la fórmula de Tweedie en $\mathbb{E}[\mathbf{x}_1 \mid \mathbf{x}_t]$, se puede establecer una relación equivalente entre el\textbf{campo de velocidad} y la\textbf{función de puntaje}:
$$u_t(\mathbf{x}_t) = \frac{\dot{\sigma}_t}{\sigma_t} \mathbf{x}_t + \left(\dot{\alpha}_t - \frac{\dot{\sigma}_t \alpha_t}{\sigma_t}\right) \cdot \frac{\mathbf{x}_t + \sigma_t^2 \nabla_{\mathbf{x}_t} \log p_t(\mathbf{x}_t)}{\alpha_t}$$

Esto significa que: dado cualquier punto de datos intermedio $\mathbf{x}_t$ y un tiempo $t$, basta con conocer la función de puntaje para calcular el campo de velocidad, y viceversa.
\end{importantbox}

\subsection{Relación entre el campo de velocidad y la predicción de ruido}

Utilizando la relación conocida entre la función de puntaje y la predicción de ruido $\nabla_{\mathbf{x}_t} \log p_t(\mathbf{x}_t) = -\frac{\boldsymbol{\epsilon}}{\sigma_t}$, podemos establecer aún más la relación equivalente entre el campo de velocidad y la red de predicción de ruido:

$$u_t(\mathbf{x}_t) = \frac{\dot{\alpha}_t}{\alpha_t} \mathbf{x}_t - \left(\dot{\alpha}_t \frac{\sigma_t}{\alpha_t} - \dot{\sigma}_t\right) \boldsymbol{\epsilon}_\theta(\mathbf{x}_t, t)$$

Donde $\boldsymbol{\epsilon}_\theta$ es la red entrenada de predicción de ruido en el modelo de difusión.

\begin{warningbox}{Predicción de velocidad vs.\ predicción de ruido: no es solo cambiar la cabeza de salida}
Aunque el campo de velocidad y el campo de ruido pueden transformarse exactamente entre sí matemáticamente, elegir predecir una u otra cantidad durante el\textbf{entrenamiento} afecta el comportamiento numérico de los gradientes. Las prácticas recientes indican que entrenar directamente la predicción de velocidad (v-prediction) ofrece mayor estabilidad numérica en ciertos esquemas de programación en comparación con la predicción de ruido ($\epsilon$-prediction), especialmente cuando $t$ se acerca a los extremos. Esto es una de las razones por las que SD3 y Flux optaron por v-prediction.
\end{warningbox}

\subsection{Consistencia con el PF-ODE de flujo de probabilidad}

El ponente señala además que el campo vectorial derivado desde Flow Matching es completamente consistente con el Probability Flow ODE (PF-ODE) derivado desde el marco SDE. Recuerde la forma general del PF-ODE:

$$d\mathbf{x} = \left[f(\mathbf{x}, t) - \frac{1}{2}g(t)^2 \nabla_\mathbf{x} \log p_t(\mathbf{x})\right] dt$$

Donde $f(\mathbf{x}, t)$ y $g(t)$ son el término de deriva y el coeficiente de difusión del SDE, respectivamente. Sus relaciones con $\alpha_t, \sigma_t$ son:
$$f(\mathbf{x}, t) = \frac{\dot{\alpha}_t}{\alpha_t} \mathbf{x}, \quad g(t)^2 = \frac{d\sigma_t^2}{dt} - 2\frac{\dot{\alpha}_t}{\alpha_t}\sigma_t^2$$

Al sustituir estos valores en el PF-ODE, se puede verificar que su resultado es\textbf{completamente idéntico} al campo vectorial derivado mediante Flow Matching.

\begin{importantbox}{Unificación de tres perspectivas equivalentes}
Hasta aquí, hemos establecido una relación de equivalencia completa entre las siguientes tres perspectivas (bajo un camino condicional gaussiano lineal):
\begin{enumerate}
    \item \textbf{Perspectiva del modelo de difusión}: aprender la función de puntaje $\nabla_\mathbf{x} \log p_t(\mathbf{x})$ o el ruido $\boldsymbol{\epsilon}$
    \item \textbf{Perspectiva de Flow Matching}: aprender el campo vectorial/velocidad $u_t(\mathbf{x})$
    \item \textbf{Perspectiva del PF-ODE}: obtener una trayectoria de muestreo determinista mediante la conversión de SDE a ODE
\end{enumerate}
Las tres son transformables exactamente entre sí en términos matemáticos; las diferencias radican en qué cantidad se predice durante el entrenamiento y qué estrategia de resolución se utiliza al momento del muestreo.
\end{importantbox}

\subsection{Resumen de la sección}

Esta sección ha establecido una relación matemática exacta de equivalencia entre los modelos de difusión y los modelos de flujo. La conclusión central es que, bajo un mapeo de flujo condicional gaussiano lineal, el modelo de difusión es un caso especial del modelo de flujo. El campo de velocidad, la función de puntaje y la predicción de ruido pueden intercambiarse libremente. Esta relación de equivalencia no solo tiene significado teórico, sino que implica en la práctica que: un modelo de difusión preentrenado puede ser directamente ``reinterpretado'' como un modelo de flujo para su uso.

